\title{Large Language Models Can Automatically Engineer Features for Few-Shot Tabular Learning}

\begin{abstract}
Large Language Models (LLMs) have shown exceptional proficiency in solving complex and novel reasoning tasks, positioning them as promising tools for tabular data analysis, a necessity in numerous practical scenarios. This study introduces an innovative in-context learning paradigm, \textsf{FeatLLM}, leveraging LLMs as feature creators to construct an input dataset that is optimally tailored for predictive modeling on tabular data. The synthesized features are subsequently utilized by a straightforward downstream machine learning algorithm—such as linear regression—to estimate class probabilities, thus enabling highly effective few-shot learning. Importantly, the \textsf{FeatLLM} architecture relies solely on the simple prediction model alongside the extracted features during inference. In contrast to prevailing LLM-driven methods, \textsf{FeatLLM} eliminates the requirement of querying the LLM per sample at inference, significantly enhancing efficiency. Additionally, it only necessitates API-level interaction with LLMs and bypasses constraints associated with prompt size. Empirical results on diverse tabular benchmarks covering multiple domains reveal that \textsf{FeatLLM} produces superior quality decision rules, outperforming alternatives like TabLLM and STUNT by a notable margin (average improvement of 10\%).

\end{abstract}

\section{Introduction}
The advancement of Large Language Models (LLMs) has led to their notable capacity for producing relevant and precise answers to novel tasks, often without specialized fine-tuning for each task~\cite{creswell2022selection,imani2023mathprompter,taylor2022galactica}. Through training on vast amounts of textual data, LLMs acquire extensive background knowledge that can serve as a substitute for expert insight in a range of contexts~\cite{Genomes2021,singhal2023large,wilcox2003role}. This enables them to streamline and automate various processes, from analyzing dialogues~\cite{finch2023leveraging} to repairing programs~\cite{fan2023automated}. Crucially, when prompted with a brief description and several sample instances, LLMs are capable of contextualizing tasks—identifying pertinent features and examples for the problem at hand~\cite{brown2020language,zhao2023survey}. These abilities underscore LLMs’ competence in data interpretation and position them as potential expert feature engineers.

Expanding upon these capabilities, researchers have started employing LLMs’ intrinsic reasoning and background knowledge in tabular learning tasks. For example, knowledge like recognizing heightened disease risks among older populations can directly aid in clinical prediction scenarios. In practice, emerging approaches articulate tasks and feature specifications in natural language, transform tabular information into serial form, and then present it to LLMs~\cite{dinh2022lift,wang2023anypredict}. Subsequent processing can involve either in-context learning techniques~\cite{nam2023semi} or methods that adjust parameters with minimal overhead~\cite{dinh2022lift,hegselmann2023tabllm}. Leveraging LLMs’ prior knowledge has yielded advantages, notably exceeding traditional tabular learning performance—especially where only a few examples are available~\cite{hegselmann2023tabllm}.

Existing approaches for leveraging large language models (LLMs) in tabular data learning face several significant hurdles. For end-to-end prediction tasks, each individual sample necessitates at least a single inference with the LLM, resulting in considerable computational overhead. Additionally, attaining high predictive accuracy often depends on fine-tuning the LLM, an impractical requirement for models that lack open-source training infrastructures or available fine-tuning APIs—an issue accentuated by the prevalence of high-performing LLMs accessible only through inference APIs~\cite{anil2023palm,openai2023gpt4}. Moreover, most current techniques exhibit poor scalability when confronted with extensive input prompts~\cite{liu2023lost}; as the feature count in tabular datasets increases and corresponding text representations lengthen, the feasibility and efficacy of the methods diminish. Collectively, these limitations fundamentally restrict the practical deployment of LLM-based tabular learning solutions.

In contrast, our methodology reconceptualizes the use of LLMs—not as direct prediction engines, but as tools for generating informative features. To achieve this, we introduce the \textsf{FeatLLM} framework, which employs in-context learning to discern the implicit ``criteria'' that guide LLM decisions. As an illustrative example, consider a disease classification scenario: the LLM is capable of articulating rules, grounded in feature values, that define the presence of the disease. These decision-making criteria are then transformed into new features that supplant the original attributes for subsequent task-solving. This paradigm shift assigns LLMs the task of feature extraction, enabling the construction of lightweight downstream models and accelerating inference time, as complex models become unnecessary once these features are established. By harnessing the in-context learning capabilities of LLMs—incorporating exemplar cases within prompts—we can generate meaningful features without explicit model training. Furthermore, adopting ensemble-inspired methods~\cite{breiman2001random}, such as feature bagging, helps address the issue of unwieldy prompt lengths.

\textsf{FeatLLM} organizes the input prompt for feature creation into two distinct phases: (1) comprehending the underlying problem to determine how features relate to target variables; and (2) formulating explicit classification rules using insights from both the previous phase and a limited set of example instances. This methodical reasoning strategy guides LLMs to disregard features lacking informative value during rule formulation. The resulting rules are executed on the data, interpreted as code, producing a binary feature to indicate each sample’s conformity with the extracted rules. Subsequently, a simple model is employed to predict class probabilities utilizing these engineered features. Stability is further enhanced through ensemble techniques that iteratively generate features with feature bagging. By judiciously managing the quantities of features and samples used in bagging, issues arising from overly large prompts are alleviated. \textsf{FeatLLM} avoids any need for additional training and exhibits minimal inference overhead, making LLM-based applications practical across diverse settings.

Our evaluation of \textsf{FeatLLM} spans 13 tabular datasets under low-shot conditions, demonstrating its reliable and superior performance. We find that LLMs efficiently exploit both prior knowledge and information derived from data, leading to the extraction of robust features. Compared to recent few-shot learning approaches, our framework achieves better performance in various scenarios. Code is provided at the anonymous GitHub repository:~\texttt{https://github.com/Sungwon-Han/FeatLLM}.

\section{Related Work}

\subsection{Few-Shot Learning with Tabular Data}

Recent progress in few-shot learning has made it possible to extract broadly applicable representations from datasets, even when labeled examples are scarce~\cite{chen2018closer}. While the initial applications were concentrated on image data, the methodology has shown robustness and adaptability, extending its effectiveness to diverse modalities such as text, audio, and notably, tabular datasets~\cite{majumder2022few,nam2023stunt,schick2021exploiting}. The challenge inherent in learning with limited annotated samples lies in the potential to model coincidental relationships that lack meaningful correlation~\cite{bartlett2021deep}. To counteract this, contemporary strategies have leveraged supplemental information sources or incorporated domain-relevant inductive biases via prior knowledge. One common avenue is semi-supervised learning, which exploits a wealth of available unlabeled data through suitable augmentation methods~\cite{ucar2021subtab,yoon2020vime}. Alternatively, unsupervised pre-training—independent of particular downstream tasks—has been utilized for effective model initialization~\cite{somepalli2021saint}. These unsupervised techniques often involve tasks such as deliberately masking or corrupting data, followed by optimization for reconstruction~\cite{arik2021tabnet,majmundar2022met}, or deploying data augmentation alongside contrastive learning paradigms~\cite{bahri2022scarf,chen2020simple}. Nonetheless, due to the intrinsic diversity within tabular datasets, devising universally appropriate augmentations remains an obstacle, and such approaches have shown limited efficacy in the few-shot regime~\cite{nam2023stunt}.

In response, recent research has adopted a transfer learning framework, training models across a wide spectrum of tabular datasets, not limited to those strictly linked to the target task~\cite{zhang2023generative,levin2022transfer,zhu2023xtab}. As an illustrative example, TransTab~\cite{wang2022transtab} employs transformer architectures to capture semantic associations among columns based on column names in addition to their values. The representations cultivated from an array of tables and column metadata enable both zero-shot and few-shot inference for unseen tabular tasks. Conversely, TabPFN~\cite{hollmann2022tabpfn} synthetically mixes distribution types to pre-train a Bayesian neural network; following pre-training, both training and testing samples from the target task are used directly for inference, without additional optimization. The incorporation of prior knowledge derived from diverse datasets provides improvements over conventional tree-based techniques and proves advantageous in few-shot learning scenarios.

\subsection{Language-Interfaced Tabular Learning}

A distinct direction for leveraging prior knowledge in tabular learning involves integration with large language models (LLMs)~\cite{anil2023palm,openai2023gpt4}. These LLMs, trained on extensive, heterogeneous textual datasets, possess considerable generalization capabilities for tasks not encountered during training, validating their applicability across numerous fields~\cite{brown2020language}. Recent investigations have explored the potential of LLMs by utilizing their embedded prior knowledge for tabular problem-solving. Approaches commonly entail converting tabular information into a textual format and supplying these as input prompts for LLMs~\cite{dinh2022lift,hegselmann2023tabllm,wang2023anypredict}. By appending tabular data, detailed feature explanations, and task summaries in the prompt, LLMs can deduce associations between features and tasks to produce inference results. The research landscape encompasses approaches such as employing parameter-efficient fine-tuning mechanisms like LoRA~\cite{hu2021lora} or IA3~\cite{liu2022few}, and using in-context learning through the inclusion of illustrative few-shot examples within prompts, obviating the need for additional training~\cite{dinh2022lift,hegselmann2023tabllm,nam2023semi,slack2023tablet}.

Although prior approaches have demonstrated success in enhancing few-shot tabular learning, their reliance on running inference entirely through the LLM imposes obstacles, particularly in industrial settings. This work seeks to move away from the traditional paradigm of processing predictions end-to-end in a black-box LLM, instead steering the LLM toward the explicit extraction of conditions or rules corresponding to each label. \textsf{FeatLLM} converts these inferred rules into executable program code that generates new features, thereby allowing downstream prediction without engaging the LLM for inference. Leveraging in-context learning, the system extracts rules by utilizing existing domain knowledge as well as information gleaned from limited sample data.

\section{Methods}
The primary function of \textsf{FeatLLM} is to distill ``rules", or conditional statements pertinent to each target class, from a small set of example inputs, illustrated in Figure~\ref{fig:main_model}. Using an LLM, these rules are generated and then translated into binary features for each sample, indicating compliance with each rule. These binary indicators are subsequently fed into a dot product with non-negative, trainable weights to compute the class probabilities. As training progresses, the model discovers which features are most relevant and influential in determining class likelihoods (see Section~\ref{Sec:training}). The process is iterated over multiple bootstrapped subsets (bagging), after which the models are combined through an ensemble mechanism. Detailed explanations of the problem formulation and method stages follow.

\vspace*{-0.12in}
\paragraph{Problem formulation.} Consider a tabular dataset containing $N$ labeled instances, $\mathcal{D} = \{(\mathbf{x}^i, \mathbf{y}^i)\}_{i=1}^{N}$. Here, each input vector $\mathbf{x}^i$ has $d$ dimensions, representing the $d$ distinct input features, while $\mathcal{D}$ is accompanied by a set of feature names $F = \{f_j\}_{j=1}^{d}$—each specified in natural language and with definitions provided elsewhere. Within the classification framework, the label $\mathbf{y}^i$ refers to an assigned class from a predefined collection. In $k$-shot learning settings, the model is trained on only $k$ labeled examples (where $k < N$), which are sampled at random.

\subsection{Prompt Design for Extracting Rules}
\label{Sec:prompt_design}
To enhance the LLM's ability to derive rules through a sound reasoning trajectory, we structure the prompt to reflect expert-level analytical strategies typical of tabular learning scenarios. As depicted in Fig.~\ref{fig:prompt_for_rule} and detailed in Appendix~\ref{sec:example_rule}, our approach organizes the prompt into three key segments:

\vspace*{-0.12in}
\paragraph{Basic information description.} This section supplies all necessary context to tackle the problem at hand (illustrated by the orange-highlighted text in Fig.~\ref{fig:prompt_for_rule}). Included are the task description, associated labeling details, specifications for each feature, and relevant demonstration examples. The task itself is described in the form of a query (such as, ``Does this patient's myocardial infarction complications data indicate chronic heart failure? Yes or no?". Additional task description examples are available in Appendix~\ref{sec:task_desc}). Feature descriptions indicate the type of value involved, and may also provide their definitions. If the feature is categorical, potential category examples are listed. For demonstration purposes, a handful of training examples—alongside their correct labels—are formatted as text. This serialization follows the conventions established in earlier works~\cite{dinh2022lift,hegselmann2023tabllm}:

\begin{align}
&\text{Serialize}(\mathbf{x}^i,\mathbf{y}^i, F) = \nonumber \\ 
&``\ f_1\ \text{is}\ \mathbf{x}^i_1.\ ... \ f_d\  \text{is}\  \mathbf{x}^i_d.\ \ \text{Answer:}\  \mathbf{y}^i\ ”
\end{align}

\vspace*{-0.12in}
\paragraph{Reasoning instruction.} To strengthen the reasoning abilities of the LLM, we deliver tailored reasoning instructions (illustrated as blue text in Fig.~\ref{fig:prompt_for_rule}). These instructions commence with an opening sentence that draws upon the chain-of-thought paradigm~\cite{wei2022chain}. We further delineate the inference workflow into two distinct stages. Initially, the LLM is prompted to hypothesize about causal links or general tendencies between the input features and the overarching task description. This phase excludes any reliance on example demonstrations and instead leverages the LLM's accumulated general or commonsense knowledge, maximizing its prior understanding of the subject matter. In the subsequent stage, the model synthesizes knowledge from the first phase alongside the provided demonstration examples to generate explicit rules pertaining to each class. Structuring reasoning in two steps discourages the model from targeting irrelevant features, thus reducing spurious correlations and enabling prioritization of key features. To manage the breadth of extracted rules—avoiding both underfitting and unnecessarily verbose prompts—we establish a constant quota for rule extraction (namely, 10 per class)\footnote{A detailed analysis regarding rule quantity is available in Section~\ref{sec:analysis}.}. Moreover, in the rule generation process, we instruct the LLM to account for each feature's type listed in the foundational information description, thereby preventing mismatched rule types.

\vspace*{-0.12in}
\paragraph{Response instruction.} To streamline the interpretation and subsequent application of the rules the LLM derives, we explicitly guide response structuring using dedicated instructions (highlighted as yellow text in Fig.~\ref{fig:prompt_for_rule}). The prompt specifies both the expected template for each answer class (i.e., response format) and the structural requirements for individual rules (i.e., rule format). These detailed guidelines enable the LLM to align its responses with the specific feature types present. In the absence of further directives, the LLM possesses the flexibility to logically combine rules through operators such as AND or OR. An illustrative outcome can be found in Appendix~\ref{sec:outcome-rule}.

\subsection{Inferring Class Likelihood via Rules}
\label{Sec:training}

\paragraph{Parsing rules for feature generation.} The rules derived in the prior stage serve as the foundation for constructing novel binary features. For each class, these features reflect whether a sample adheres to the pertinent rules—resulting in a value of either '0' or '1'. Such binary features can subsequently inform the estimation of a sample's propensity to belong to a particular class. Nevertheless, because rules originating from the LLM are articulated in natural language, it becomes necessary to implement a parsing procedure to automate feature extraction. While efforts are made to enforce response and rule formatting for improved parseability, outputs from the LLM may still manifest inconsistent syntactic variations~\cite{kovavc2023large}, such as replacing symbols like ``$>$” or ``$<$” with verbal equivalents (e.g., ``greater than'' or ``less than''). These inconsistencies complicate automatic parsing.

To overcome the intricacies associated with parsing imprecisely formatted text, instead of resorting to elaborate parsing code, we take advantage of the LLM’s capabilities. Since the LLM excels at grasping semantic intent amid textual variability and can synthesize straightforward code when supplied with instructions—leveraging its training on code-related data—we provide a prompt that encompasses function names, descriptions of inputs and outputs, and the extracted rules. This prompt is then submitted to the LLM. Illustrations of these prompts as well as sample outcomes can be found in Figures~\ref{fig:prompt_for_parsing} and ~\ref{sec:example_parsing}-\ref{sec:example_parse_outcome} in the Appendix. The resulting code is subsequently executed via Python’s \texttt{exec()} function to facilitate the conversion of data.

\vspace*{-0.12in}
\paragraph{Inferring class likelihood.} With the new binary features in place—each corresponding to rules for a specific class—a straightforward technique for quantifying a sample’s likelihood for belonging to any class is to tally the satisfied rules (i.e., aggregate the binary feature values for that class). Nevertheless, because rules and features may hold unequal predictive value, their relative importance should be learned from available training data. We address this by employing standard machine learning (ML) methods to model the relevance of each binary feature, applied for every class. As an instance, a bias-free linear model can be used. If the binary feature vector created from sample $\mathbf{x}^i$ for class $k$ is denoted as $\mathbf{z}_k^i \in \{0, 1\}^R$, where $R$ signifies the rule count per class and $c$ the total class count, the class probability $\mathbf{p}^i$ can be calculated as:

\begin{align}
\text{logit}_k^i &= \max(\mathbf{w}_k, 0) \cdot \mathbf{z}_k^i, \nonumber \\
\mathbf{p}^i &= \text{Softmax}({[\text{logit}_{1}^i, ..., \text{logit}_{c}^i]}). \label{eq:infer_prob}
\end{align}

In this context, $\mathbf{w}_k$ denotes the adjustable weights of the linear model, which serve to reflect the relevance of each individual feature. By constraining these weights to reside within a non-negative space, we guarantee that features derived from the LLM do not diminish the probability assigned to a given class during training. Subsequently, the resulting logit vector is transformed into class probabilities via the application of the softmax function. The learning of $\mathbf{w}_k$ is achieved by minimizing the cross-entropy loss using a limited set of training examples.

\vspace*{-0.12in}
\paragraph{Ensembling with bagging.} In the final step, we repeat the preceding procedure several times to generate an ensemble of inference models, ultimately aggregating their outputs for prediction. For $T$ repetitions, the corresponding set of learned weights is represented as $\{\mathbf{w}[t]\}_{t=1}^T$. The final class prediction is based on the average probability for each class $\mathbf{p}^i$, where each is determined using a specific weight $\mathbf{w}[t]$, as described in Eq.~\ref{eq:infer_prob}.

To increase diversity during rule generation for the ensemble, the LLM's inference operates at an elevated temperature. In addition, the order of few-shot examples is randomized, reflecting the insight that varying example sequencing affects LLM outputs~\cite{lu2022fantastically}\footnote{Appendix~\ref{sec:diversity} provides further analysis of rule diversity}. Bagging is utilized by selecting random subsets of features or training instances for each run, especially advantageous with greater numbers of features or samples. This ensembling methodology confers two primary benefits. First, it compensates for erroneous rules generated by the LLM in certain runs, as accurate rules from other iterations bolster the overall robustness. This aligns with the principle of LLM self-consistency~\cite{wang2022self}, whereby rules persistently produced across trials are likelier to be correct. Second, it effectively mitigates constraints imposed by the LLM's input prompt length; bagging allows for selective input of features or instances, facilitating resilient model performance even with large datasets. Whereas any single trial's rule set may fail to account for all feature relationships, the aggregate of weak learners across multiple iterations diminishes the drawbacks of incomplete coverage in individual runs.

\section{Experiments}
We assess the effectiveness of \textsf{FeatLLM} across a range of tabular datasets and perform thorough investigations to better understand the mechanisms of our proposed method. For brevity, supplementary analyses—including LLM variants, qualitative evaluations, and further experimental results—are available in the Appendix. 

\subsection{Performance Evaluation}
\paragraph{Datasets.}
Our study employs 13 datasets suitable for either binary or multiclass classification tasks: (1) Adult~\cite{asuncion2007uci}, used to predict whether a person's annual income exceeds \$50,000; (2) Bank~\cite{moro2014data}, which forecasts customer subscription to term deposits; (3) Blood~\cite{yeh2009knowledge}, aimed at identifying donors likely to return; (4) Car~\cite{kadra2021well}, for evaluating car quality; (5) Communities~\cite{misc_communities_and_crime_183}, which estimates regional crime rates; (6) Credit-g~\cite{kadra2021well}, for assessing credit risk of individuals; (7) Diabetes\footnote{\texttt{kaggle.com/datasets/uciml/pima-indians-diabetes-database}}, to determine diabetes diagnosis; (8) Heart\footnote{\texttt{kaggle.com/datasets/fedesoriano/heart-failure-prediction}}, which predicts the likelihood of coronary artery disease; and (9) Myocardial~\cite{misc_myocardial_infarction_complications_579}, concerning chronic heart failure occurrence. 

We expand this collection with contemporary and synthetic datasets which are rarely addressed in prior literature and excluded from LLM pre-training: (10) Cultivars, tasked with predicting soybean cultivar grain yield\footnote{\texttt{archive.ics.uci.edu/dataset/913}}; (11) NHANES, focused on classifying individuals by age group\footnote{\texttt{archive.ics.uci.edu/dataset/887}}; (12) Sequence-type, which categorizes sequences as arithmetic, geometric, Fibonacci, or Collatz; (13) Solution-mix, identifying whether a mixture of four solutions exceeds 0.5 concentration. The first two among these datasets were deposited into the UCI Machine Learning Repository after September 2023, ensuring their absence from the training data of the LLM API (gpt-3.5-turbo-0613) underpinning our model. The last two datasets are artificially generated, precluding any prior exposure by the LLM API and facilitating unbiased evaluation.

All datasets exhibit varying degrees of scale and difficulty, detailed comprehensively in Table~\ref{Tab:basic-desc}.
Each data source features well-documented attribute labels and descriptions. Some datasets, such as `Communities' and `Myocardial', comprise over 100 features, posing practical difficulties for LLMs given context window limitations. 

\vspace*{-0.12in}
\paragraph{Baselines.}
For benchmarking, we compare our approach against nine distinct baselines. The first three represent classic tabular modeling: (1) Logistic regression (LogReg), (2) XGBoost~\cite{chen2016xgboost}, and (3) RandomForest~\cite{ho1995random}. The next two assume access to additional unlabeled data: (4) SCARF~\cite{bahri2022scarf}, and (5) STUNT~\cite{nam2023stunt}. Note that gathering substantial unlabeled datasets can be difficult in practical scenarios. Another recent comparison point is (6) TabPFN~\cite{hollmann2022tabpfn}, which leverages synthetic data with varied distributions to pre-train models. The last three baselines are based on LLM methodologies: (7) In-context learning (In-context)~\cite{wei2022emergent}, (8) TABLET~\cite{slack2023tablet}, and (9) TabLLM~\cite{hegselmann2023tabllm}. In-context learning directly incorporates limited training instances into prompts, without model adjustments. TABLET enhances prompt quality through rule-based hints generated by an external classifier. TabLLM relies on parameter-efficient approaches like IA3~\cite{liu2022few} to fine-tune LLMs using a small number of training samples.

\vspace*{-0.12in}
\paragraph{Implementation details.}
\textsf{FeatLLM} leverages GPT-3.5 as its LLM foundation but has been constructed to function independently of any specific LLM architecture (refer to Appendix~\ref{sec:backbones} for experiments using alternative LLMs). For inference, the temperature parameter is configured at 0.5, and the top-p parameter adopts the API's standard value of 1. The ensemble count and rule extraction limit are fixed at 20 and 10, respectively. Insights on how these hyper-parameters affect performance are shown in Figure~\ref{fig:hyperparameter2} and detailed in Appendix~\ref{sec:hyper-parameter}. Linear model training uses the Adam optimizer, set with a learning rate of 0.01, spanning 200 training epochs. To select the optimal epoch, $k$-fold cross-validation is applied. When evaluating baselines, in-context learning leverages GPT-3.5, while TabLLM utilizes T0~\cite{sanh2022multitask}, matching the methodology specified in the original source. It is important to emphasize that TabLLM confines LLM selection to those models with accessible public checkpoints, rendering GPT-3.5 unavailable for this benchmark. Conventional machine learning algorithms such as LogReg, XGBoost, and RandomForest have their parameters chosen via Grid-search integrated with $k$-fold cross-validation. Here, the value of $k$ is either 2 or 4 to ensure that the training partition contains at least one representative from each class. Comprehensive implementation details can be found in Appendix~\ref{sec:baselines}.

\vspace*{-0.12in}
\paragraph{Results.}
Tables~\ref{tab:main1} and \ref{tab:main2} provide comparative results for different models across a diverse set of datasets. To ensure reliable evaluation, each experiment is conducted three times with varied random splits for training and testing sets; we report both the mean and standard deviation of AUC scores. \textsf{FeatLLM} demonstrates robust performance, often emerging as the leading approach or placing second relative to other baselines. Remarkably, our method matches the effectiveness of traditional tabular baselines using only 4 shots, whereas those baselines need 32 shots for similar results (refer to Fig.~\ref{fig:compares}). Although STUNT and TabPFN displayed notable gains on select datasets, their improvements over conventional machine learning approaches were modest on average. Furthermore, LLM-based solutions—including in-context learning, TABLET, and TabLLM—were unable to make inferences on datasets with high feature dimensionality. In contrast, our strategy, which integrates bagging and ensemble methods, remains capable of delivering strong performance even as feature count grows. Additionally, it is worth emphasizing that the bagging and ensemble framework we propose can be incorporated into other LLM-based methods; nonetheless, this would double the inference time per sample, thereby greatly increasing computational demands and reducing feasibility.

\subsection{Analysis \& Discussion}
\label{sec:analysis}
\paragraph{Ablation study.} To discern the role of each major component in our architecture, we conduct ablation studies, examining the impact of excluding: (1) \textsf{-Tuning}, by removing the linear model's weight adjustment and simply aggregating binary features to compute the logit; (2) \textsf{-Ensemble}, by omitting the ensemble step; (3) \textsf{-Description}, by excluding feature descriptions; and (4) \textsf{-Reasoning}, by skipping the Step-1 reasoning instruction during rule generation.

Table~\ref{tab:ablation} details the average AUC change attributable to each ablation across all evaluated datasets. Omitting any single component consistently causes a decrease in model performance. Of note, the advantages of weight tuning are magnified as the number of shots grows, suggesting that greater data availability allows for more precise rule weighting and heightens the value of tuning. Conversely, removing feature descriptions and reasoning instructions causes more pronounced performance drops in low-shot settings, which implies that leveraging LLMs' prior knowledge becomes especially crucial when limited examples are available. Finally, employing our proposed ensemble methodology reliably boosts results under every experimental condition.

\vspace*{-0.12in}
\paragraph{Training \& inference time comparison.} We analyze and report the training and inference durations among multiple methods. The experiments utilize the Adult dataset, specifically using a training subset with 16 shots. All experiments are carried out on a single A100 GPU as the standard setting, except for TabLLM, which operates across four A100 GPUs to facilitate model parallelism. Table~\ref{tab:cost} details these runtime comparisons. Remarkably, \textsf{FeatLLM} achieves inference time that is nearly as efficient as traditional approaches. Conversely, alternative LLM-driven methods—including In-context learning, TABLET, and TabLLM—demonstrate substantially greater inference latency. With respect to training duration, \textsf{FeatLLM} maintains parity with other LLM-centric techniques, necessitating only 30 API requests. This does not introduce meaningful overhead to the cost and, importantly, can be parallelized.

\vspace*{-0.12in}
\paragraph{Quality of features generated by \textsf{FeatLLM}.}
We present a straightforward analysis to evaluate the informativeness of rule-based features produced by \textsf{FeatLLM} for real-world tasks. In particular, we assess the average AUROC for newly created features relative to the target labels, followed by calculating the proportion of features attaining AUROC scores exceeding 0.5 within each dataset. Table~\ref{tab:quality} compiles these findings. The results demonstrate that most of the generated rules align well with the target attribute and yield meaningful information pertinent to solving the task (see the ``Before tuning'' column). Furthermore, during the tuning phase of the linear model using training data, rules exhibiting low utility (i.e., those associated with weights beneath the specified threshold) are systematically excluded (refer to the ``After tuning'' column). Here, the threshold is fixed at 0.1, in accordance with the global weight histogram.

\vspace*{-0.12in}
\paragraph{Effect of adding spurious correlations.} To evaluate the effectiveness of different approaches in mitigating the impact of spurious correlations, especially in low-shot scenarios, we performed an investigation. In this experiment, we used the Heart dataset and appended random columns from the Adult dataset, which hold no relevance to the classification objective of the Heart dataset. We then examined how model performances respond as the number of extraneous columns increases. For comparability, all methods were given 8 labeled examples, without inclusion of any unlabeled instances; this omits methods such as SCARF and STUNT from our evaluation. 

Figure~\ref{fig:spurious} depicts how model accuracy is affected as spurious correlations are introduced. Relative to baseline methods, \textsf{FeatLLM} exhibited minimal decline in performance when faced with these irrelevant columns. This outcome appears attributable to the framework’s reliance on prior task-feature associations during rule extraction, which inhibits the creation of rules predicated on irrelevant features, thereby promoting resilience against noise. Our analysis confirms that rules generated from noisy columns comprised only around 1-2\% of the rule set. Nonetheless, it is also evident that \textsf{FeatLLM} remains susceptible to learning false associations and misinterpreting causal links between task and features when the added columns are from data with underlying similarities (see Appendix~\ref{sec:spurious-appendix}).

\vspace*{-0.12in}
\paragraph{Hyper-parameter analysis}
\textsf{FeatLLM} is characterized by two primary hyper-parameters: the ensemble count and the number of rules extracted per class during each LLM invocation. To understand how these settings affect model outcomes, we performed a series of tests. Our findings indicate that increasing the ensemble number generally leads to enhanced performance; however, this improvement plateaus once the number reaches approximately 20 (refer to Fig.~\ref{fig:appendix-hyperparameter1} in Appendix). Regarding the rule quantity, no significant difference was observed among the choices, but setting this value at 10 was deemed optimal, balancing prompt complexity and predictive accuracy (see Fig.~\ref{fig:hyperparameter2}). We hypothesize that a higher rule count expands input dimensionality, contributing additional information, though it also tends to induce overfitting, especially in scenarios with limited samples.

\section{Conclusion}
We introduce \textsf{FeatLLM}, which establishes new benchmarks for LLM-driven few-shot learning on tabular datasets. Unlike conventional approaches that merely employ LLMs for individual sample prediction, \textsf{FeatLLM} emphasizes the derivation of criteria for prediction, mimicking feature engineering practices. By leveraging both rule extraction and ensemble bagging techniques, \textsf{FeatLLM} maintains low inference demands, requires only access to LLM APIs without the necessity for model training, and surpasses existing methods in handling feature dimension constraints. As a result, \textsf{FeatLLM} attains impressive prediction results and serves as a highly efficient solution for real-world tabular data tasks using LLMs.

While \textsf{FeatLLM} is specifically designed for low-shot learning situations, future work aims to extend its feature creation capabilities to datasets with larger sample sizes. This includes investigating diverse feature representations beyond rule-based schemes and enhancing interpretability, facilitating broader applicability to multiple domains.

\section{Broader Impact}
The \textsf{FeatLLM} framework leverages LLMs' prior knowledge to improve performance in tabular learning scenarios, surpassing conventional supervised learning under data constraints. By supplementing scarce training data with pretrained knowledge, our work addresses the overfitting issue common in low-resource settings. \textsf{FeatLLM} is especially useful for practitioners in sectors such as finance and healthcare, where labeling costs are prohibitive or collecting labels presents significant challenges, and domain-relevant LLM expertise can greatly augment learning (as LLMs typically encode substantial information from finance or healthcare-related sources). For broader societal uptake, a promising direction is to investigate the consequences of LLMs’ embedded biases and misinformation, as these prior beliefs are likely to influence, and possibly outweigh, the impact of labeled training data in determining model outputs.

\end{document}